\ifdefined\gkver\else\def\gkver{student}\fi
\documentclass[\gkver]{gaokaozhenti}
\begin{document}

\chapter{2000年北京}

\begin{enumerate}
\item
%% number: 1
%% paperName: 2000年北京春考第 1 题 4 分
%% typeId: 2
%% type: 多选题
%% chapter: 光学
%% point: 光电效应
%% method: 无
%% score: 4
%% degree: 750
%% duplicateId: 0
%% body:
某种金属在单色光照射下发射出光电子，这些光电子的最大初动能 \xzanswer{BD}

\fourchoices[answer=BD]
{随照射光强度的增大而增大}
{随照射光频率的增大而增大}
{随照射光波长的增大而增大}
{与照射光的照射时间无关}

%% answer: BD
%% memo:
\memoanswer{本题原卷及材料中未提供详解，待补充。}

\item
%% number: 2
%% paperName: 2000年北京春考第 2 题 4 分
%% typeId: 1
%% type: 单选题
%% chapter: 曲线运动
%% point: 平抛运动
%% method: 无
%% score: 4
%% degree: 650
%% duplicateId: 0
%% body:
做平抛运动的物体，每秒的速度增量总是 \xzanswer{A}

\fourchoices[answer=A]
{大小相等，方向相同}
{大小不等，方向不同}
{大小相等，方向不同}
{大小不等，方向相同}

%% answer: A
%% memo:
\memoanswer{本题原卷及材料中未提供详解，待补充。}

\item
%% number: 3
%% paperName: 2000年北京春考第 3 题 4 分
%% typeId: 2
%% type: 多选题
%% chapter: 电场
%% point: 电容
%% method: 无
%% score: 4
%% degree: 650
%% duplicateId: 0
%% body:
图示是一个由电池、电阻 $R$ 与平行板电容器组成的串联电路。在增大电容器两极板间距离的过程中，\xzanswer{BC}
\onepicture[w=5cm]{figs/03.png}

\fourchoices[answer=BC]
{电阻 $R$ 中没有电流}
{电容器的电容变小}
{电阻 $R$ 中有从 a 流向 b 的电流}
{电阻 $R$ 中有从 b 流向 a 的电流}

%% answer: BC
%% memo:
\memoanswer{本题原卷及材料中未提供详解，待补充。}

\item
%% number: 4
%% paperName: 2000年北京春考第 4 题 4 分
%% typeId: 1
%% type: 单选题
%% chapter: 磁场
%% point: 带电粒子在磁场中的运动
%% method: 无
%% score: 4
%% degree: 650
%% duplicateId: 0
%% body:
一带电质点在匀强磁场中作圆周运动，现给定了磁场的磁感应强度，带电质点的质量和电量。若用 $v$ 表示带电质点运动的速率，$R$ 表示其轨道半径，则带电质点运动的周期 \xzanswer{B}

\fourchoices[answer=B]
{与 $v$ 有关，与 $R$ 有关}
{与 $v$ 无关，与 $R$ 无关}
{与 $v$ 有关，与 $R$ 无关}
{与 $v$ 无关，与 $R$ 有关}

%% answer: B
%% memo:
\memoanswer{本题原卷及材料中未提供详解，待补充。}

\item
%% number: 5
%% paperName: 2000年北京春考第 5 题 4 分
%% typeId: 1
%% type: 单选题
%% chapter: 原子和原子核
%% point: 玻尔模型
%% method: 无
%% score: 4
%% degree: 650
%% duplicateId: 0
%% body:
根据玻尔理论，某原子的电子从能量为 $E$ 的轨道跃迁到能量为 $E'$ 的轨道，辐射出波长为 $\lambda$ 的光。以 $h$ 表示普朗克常量，$c$ 表示真空中的光速，则 $E'$ 等于 \xzanswer{C}

\fourchoices[answer=C]
{$E - h\frac{\lambda}{c}$}
{$E + h\frac{\lambda}{c}$}
{$E - h\frac{c}{\lambda}$}
{$E + h\frac{c}{\lambda}$}

%% answer: C
%% memo:
\memoanswer{本题原卷及材料中未提供详解，待补充。}

\item
%% number: 6
%% paperName: 2000年北京春考第 6 题 4 分
%% typeId: 2
%% type: 多选题
%% chapter: 电磁感应
%% point: 远距离输电
%% method: 无
%% score: 4
%% degree: 750
%% duplicateId: 0
%% body:
远距离输送交流电都采用高压输电，我国正在研究用比 $330\UkV$ 高得多的电压进行输电，采用高压输电的优点是 \xzanswer{AC}

\fourchoices[answer=AC]
{可节省输电线的铜材料}
{可根据需要调节交流电的频率}
{可减少输电线上的能量损失}
{可加快输电的速度}

%% answer: AC
%% memo:
\memoanswer{本题原卷及材料中未提供详解，待补充。}

\item
%% number: 7
%% paperName: 2000年北京春考第 7 题 4 分
%% typeId: 2
%% type: 多选题
%% chapter: 曲线运动
%% point: 人造地球卫星
%% method: 无
%% score: 4
%% degree: 650
%% duplicateId: 0
%% body:
可以发射一颗这样的人造地球卫星，使其圆轨道 \xzanswer{CD}

\fourchoices[answer=CD]
{与地球表面上某一纬度线（非赤道）是共面同心圆}
{与地球表面上某一经度线所决定的圆是共面同心圆}
{与地球表面上的赤道线是共面同心圆，且卫星相对地球表面是静止的}
{与地球表面上的赤道线是共面同心圆，但卫星相对地球表面是运动的}

%% answer: CD
%% memo:
\memoanswer{本题原卷及材料中未提供详解，待补充。}

\item
%% number: 8
%% paperName: 2000年北京春考第 8 题 4 分
%% typeId: 1
%% type: 单选题
%% chapter: 机械振动
%% point: 单摆
%% method: 无
%% score: 4
%% degree: 750
%% duplicateId: 0
%% body:
已知在单摆 a 完成 10 次全振动的时间内，单摆 b 完成 6 次全振动，两摆长之差为 $1.6\Um$。则两单摆摆长 $l_{a}$ 与 $l_{b}$ 分别为 \xzanswer{B}

\fourchoices[answer=B]
{$l_{a} = 2.5\Um$，$l_{b} = 0.9\Um$}
{$l_{a} = 0.9\Um$，$l_{b} = 2.5\Um$}
{$l_{a} = 2.4\Um$，$l_{b} = 4.0\Um$}
{$l_{a} = 4.0\Um$，$l_{b} = 2.4\Um$}

%% answer: B
%% memo:
\memoanswer{本题原卷及材料中未提供详解，待补充。}

\item
%% number: 9
%% paperName: 2000年北京春考第 9 题 4 分
%% typeId: 1
%% type: 单选题
%% chapter: 磁场
%% point: 磁力矩
%% method: 无
%% score: 4
%% degree: 650
%% duplicateId: 0
%% body:
一矩形线圈 $abcd$ 处在磁感应强度为 $B$ 的匀强磁场中，磁场方向与 $ab$ 垂直。当线圈以角速度 $\omega$ 绕 $ab$ 转动时，感应电动势的最大值为 $E_{1}$，线圈受到的最大磁力矩为 $M_{1}$；当以角速度 $\omega$ 绕中心轴 $OO'$ 转动时，感应电动势的最大值为 $E_{2}$，最大磁力矩为 $M_{2}$。则 $E_{1}:E_{2}$ 和 $M_{1}:M_{2}$ 分别为 \xzanswer{A}
\onepicture[w=5cm]{\includesvg{figs/09}}

\fourchoices[answer=A]
{1:1，1:1}
{1:1，1:2}
{1:2，1:1}
{1:2，1:2}

%% answer: A
%% memo:
\memoanswer{本题原卷及材料中未提供详解，待补充。}

\item
%% number: 10
%% paperName: 2000年北京春考第 10 题 4 分
%% typeId: 1
%% type: 单选题
%% chapter: 电场
%% point: 电场线
%% method: 无
%% score: 4
%% degree: 650
%% duplicateId: 0
%% body:
如图所示，P、Q 是两个电量相等的正的点电荷，它们连线的中点是 O，A、B 是中垂线上的两点，$OA < OB$，用 $E_{A}$、$E_{B}$、$\varphi_{A}$、$\varphi_{B}$ 分别表示 A、B 两点的场强和电势，则 \xzanswer{B}
\onepicture[w=4.5cm]{\includesvg{figs/10}}

\fourchoices[answer=B]
{$E_{A}$ 一定大于 $E_{B}$，$\varphi_{A}$ 一定大于 $\varphi_{B}$}
{$E_{A}$ 不一定大于 $E_{B}$，$\varphi_{A}$ 一定大于 $\varphi_{B}$}
{$E_{A}$ 一定大于 $E_{B}$，$\varphi_{A}$ 不一定大于 $\varphi_{B}$}
{$E_{A}$ 不一定大于 $E_{B}$，$\varphi_{A}$ 不一定大于 $\varphi_{B}$}

%% answer: B
%% memo:
\memoanswer{本题原卷及材料中未提供详解，待补充。}

\item
%% number: 11
%% paperName: 2000年北京春考第 11 题 4 分
%% typeId: 2
%% type: 多选题
%% chapter: 运动和力
%% point: 动量和能量
%% method: 无
%% score: 4
%% degree: 550
%% duplicateId: 0
%% body:
一轻质弹簧，上端悬挂于天花板，下端系一质量为 $M$ 的平板，处在平衡状态。一质量为 $m$ 的均匀环套在弹簧外，与平板的距离为 $h$，如图所示。让环自由下落，撞击平板。已知碰后环与板以相同的速度向下运动，使弹簧伸长。 \xzanswer{AC}
\onepicture[h=5.5cm]{figs/11.png}

\fourchoices[answer=AC]
{若碰撞时间极短，则碰撞过程中环与板的总动量守恒}
{若碰撞时间极短，则碰撞过程中环与板的总机械能守恒}
{环撞击板后，板的新的平衡位置与 $h$ 的大小无关}
{在碰后板和环一起下落的过程中，它们减少的动能等于克服弹簧力所作的功}

%% answer: AC
%% memo:
\memoanswer{本题原卷及材料中未提供详解，待补充。}

\item
%% number: 12
%% paperName: 2000年北京春考第 12 题 4 分
%% typeId: 2
%% type: 多选题
%% chapter: 机械波
%% point: 波与振动图象
%% method: 无
%% score: 4
%% degree: 450
%% duplicateId: 0
%% body:
已知平面简谐波在 $x$ 轴上传播，原点 $O$ 的振动图线如图~\ref{2000北京12a} 所示，$t$ 时刻的波形图线如图~\ref{2000北京12b} 所示，则 $t' = t + 0.5\Us$ 时刻的波形图线可能是 \xzanswer{CD}

\twopicture[labelA=2000北京12a, labelB=2000北京12b, h=3cm]
{figs/12a.png}{figs/12b.png}

\fourchoices[answer=CD, ispicture=true, h=2.2cm]
{figs/12c.png}
{figs/12d.png}
{figs/12e.png}
{figs/12f.png}

%% answer: CD
%% memo:
\memoanswer{本题原卷及材料中未提供详解，待补充。}

\item
%% number: 13
%% paperName: 2000年北京春考第 13 题 5 分
%% typeId: 3
%% type: 填空题
%% chapter: 电路
%% point: 电功电功率
%% method: 无
%% score: 5
%% degree: 850
%% duplicateId: 0
%% body:
一个蓄电池输出电压为 $12\UV$，若输出的电能为 $4.3\times 10^{6}\UJ$，则它放出的电量为 \tkanswer{$3.6\times 10^{5}$}$\UC$。

%% answer: $3.6\times 10^{5}$
\jdanswer{
$3.6\times 10^{5}\UC$
}

%% memo:
\memoanswer{
由 $W = UQ$ 可得
$Q = \frac{W}{U} = \frac{4.3 \times 10^{6}}{12}\UC \approx 3.6 \times 10^{5}\UC$。
}

\item
%% number: 14
%% paperName: 2000年北京春考第 14 题 5 分
%% typeId: 3
%% type: 填空题
%% chapter: 交流电
%% point: 变压器
%% method: 无
%% score: 5
%% degree: 550
%% duplicateId: 0
%% body:
如图所示，理想变压器的原、副线圈匝数之比为 $n_{1}:n_{2} = 4:1$，原线圈回路中的电阻 A 与副线圈回路中的负载电阻 B 的阻值相等。a、b 端加一定交流电压后，两电阻消耗的电功率之比 $P_{A}:P_{B} =$ \tkanswer{1:16}，两电阻两端电压之比 $U_{A}:U_{B} =$ \tkanswer{1:4}。
\onepicture[w=5cm, align=right]{figs/14.png}

%% answer: 1:16，1:4
\jdanswer{
$P_{A}:P_{B} = 1:16$，$U_{A}:U_{B} = 1:4$
}

%% memo:
\memoanswer{
设原线圈中的电流为 $I_{1}$，副线圈中的电流为 $I_{2}$，由变压器电流与匝数成反比可得
$\frac{I_{1}}{I_{2}} = \frac{n_{2}}{n_{1}} = \frac{1}{4}$。
电阻 A、B 阻值相等，设为 $R$，则
$\frac{P_{A}}{P_{B}} = \frac{I_{1}^{2}R}{I_{2}^{2}R} = \frac{I_{1}^{2}}{I_{2}^{2}} = \frac{1}{16}$，
$\frac{U_{A}}{U_{B}} = \frac{I_{1}R}{I_{2}R} = \frac{I_{1}}{I_{2}} = \frac{1}{4}$。
}

\item
%% number: 15
%% paperName: 2000年北京春考第 15 题 5 分
%% typeId: 3
%% type: 填空题
%% chapter: 运动和力
%% point: 变加速直线运动
%% method: 无
%% score: 5
%% degree: 650
%% duplicateId: 0
%% body:
1999 年 11 月 20 日，我国发射了“神舟号”载人飞船，次日载人舱着陆，实验获得成功。载人舱在将要着陆之前，由于空气阻力作用有一段匀速下落过程。若空气阻力与速度的平方成正比，比例系数为 $k$，载人舱的质量为 $m$，则此过程中载人舱的速度应为 \tkanswer{$\sqrt{\frac{mg}{k}}$}。

%% answer: $\sqrt{\frac{mg}{k}}$
\jdanswer{
$\sqrt{\frac{mg}{k}}$
}

%% memo:
\memoanswer{
载人舱匀速下落时，重力和空气阻力平衡，有
$mg = kv^{2}$，
解得
$v = \sqrt{\frac{mg}{k}}$。
}

\item
%% number: 16
%% paperName: 2000年北京春考第 16 题 5 分
%% typeId: 3
%% type: 填空题
%% chapter: 曲线运动
%% point: 万有引力定律
%% method: 无
%% score: 5
%% degree: 450
%% duplicateId: 0
%% body:
地核的体积约为整个地球体积的 16\%，地核的质量约为地球质量的 34\%。经估算，地核的平均密度为 \tkanswer{$1.2\times 10^{4}$}$\Ukgmc$。（结果取两位有效数字）

%% answer: $1.2\times 10^{4}$
\jdanswer{
$1.2\times 10^{4}\Ukgmc$
}

%% memo:
\memoanswer{
设地球的体积为 $V$、质量为 $M$，则地球的平均密度 $\rho_{0} = \frac{M}{V}$。
地核的体积为 $0.16V$，质量为 $0.34M$，则地核的平均密度
$\rho = \frac{0.34M}{0.16V} = 2.125\rho_{0}$。
取地球的平均密度 $\rho_{0} \approx 5.6 \times 10^{3}\Ukgmc$，可得
$\rho \approx 1.2 \times 10^{4}\Ukgmc$。
}

\item
%% number: 17
%% paperName: 2000年北京春考第 17 题 4 分
%% typeId: 4
%% type: 实验题
%% chapter: 直线运动
%% point: 螺旋测微器
%% method: 无
%% score: 4
%% degree: 950
%% duplicateId: 0
%% body:
用螺旋测微器测量一根粗细均匀的金属丝的直径，示数如图所示，读数是 \tkanswer{13.857}$\Umm$。
\onepicture[w=6cm, align=right]{figs/17.png}

%% answer: 13.857（13.858 或 13.856 均可）
\jdanswer{
$13.857\Umm$（$13.858\Umm$ 或 $13.856\Umm$ 均可）
}

%% memo:
\memoanswer{
螺旋测微器的固定刻度读数为 $13.5\Umm$，可动刻度中第 $35.7$ 格与基准线对齐，
故读数为 $13.5\Umm + 35.7 \times 0.01\Umm = 13.857\Umm$（$13.858\Umm$ 或 $13.856\Umm$ 均可）。
}

\item
%% number: 18
%% paperName: 2000年北京春考第 18 题 6 分
%% typeId: 4
%% type: 实验题
%% chapter: 电路
%% point: 多用表的使用
%% method: 无
%% score: 6
%% degree: 750
%% duplicateId: 0
%% body:
在使用多用表测电阻时，有一个步骤是将两表笔短路，调节表上某旋钮，使指针指到表盘上电阻的示数为 \tkanswer{0} 处。测量某电池的端电压时，要将红表笔接到电池的 \tkanswer{正} 极。测量电流时，要使电流由 \tkanswer{红} 表笔流入多用电表。

%% answer: 0，正，红
\jdanswer{
0，正，红
}

%% memo:
\memoanswer{
使用多用电表的欧姆挡测电阻前，应先将两表笔短接，调节欧姆调零旋钮，使指针指到表盘上电阻的示数为 0 处。
测量电池的端电压时，电流从电池正极流出，故红表笔应接电池的正极。
测量电流时，应使电流从红表笔流入多用电表。
}

\item
%% number: 19
%% paperName: 2000年北京春考第 19 题 8 分
%% typeId: 4
%% type: 实验题
%% chapter: 气体性质
%% point: 查理定律
%% method: 无
%% score: 8
%% degree: 550
%% duplicateId: 0
%% body:
利用“验证玻意耳-马略特定律”的实验装置来验证查理定律。

\begin{enumerate}
\item
为了完成这个实验，除了图中给出的器材外，还需要气压计、托盘天平、热水、凉水和 \tkanswer{温度计}。

\item
必须进行的实验步骤有：

\begin{steps}
\item
用托盘天平称出活塞和框架的质量 $M$，用气压计读出实验室中的大气压强 $p_{0}$。按图安装器材，在框架两侧挂上钩码，使注射器的下半部分位于量杯之中。往量杯中加入适量的凉水，使注射器内的空气柱位于水面之下。过几分钟后，记下钩码的质量和活塞下表面的位置。同时 \tkanswer{用温度计测出杯中水的温度}；

\item
在量杯中加些热水，过几分钟后在框架两侧加挂适当质量的钩码，使 \tkanswer{活塞下表面再恢复到步骤 1 中的位置}，记下钩码的质量。同时，\tkanswer{用温度计测量水的温度}；

\item
把步骤 2 重复 4 次。
\end{steps}

\item
可用作图法来验证查理定律是否成立，该图线的横坐标所代表的物理量及其单位是 \tkanswer{温度，$\UdC$（或 $\UK$）}，纵坐标所代表的物理量及其单位是 \tkanswer{压强，$\UPa$}。
\end{enumerate}
\onepicture[h=8cm, align=right]{figs/19.png}

%% answer: （一）温度计；（二）1．用温度计测出杯中水的温度。2．使活塞下表面再恢复到步骤 1 中的位置，同时用温度计测量水的温度。（三）温度，℃（或 K）；压强，Pa（纵、横坐标互换也对）
\jdanswer{
\begin{enumerate}
	\item
	温度计。

	\item
	用温度计测出杯中水的温度；使活塞下表面再恢复到步骤 1 中的位置，同时用温度计测量水的温度。

	\item
	温度，$\UdC$（或 $\UK$）；压强，$\UPa$（纵、横坐标互换也对）。
\end{enumerate}
}

%% memo:
\memoanswer{
（一）温度计。

（二）1．用温度计测出杯中水的温度。

2．使活塞下表面再恢复到步骤 1 中的位置，同时用温度计测量水的温度。

（三）温度，$\UdC$（或 $\UK$）；压强，$\UPa$（纵、横坐标互换也对）。
}

\item
%% number: 20
%% paperName: 2000年北京春考第 20 题 12 分
%% typeId: 6
%% type: 计算题
%% chapter: 光学
%% point: 光的折射
%% method: 无
%% score: 12
%% degree: 750
%% duplicateId: 0
%% body:
\begin{enumerate}
\item
用简明的语言表述临界角的定义。

\item
玻璃和空气相接触，试画出入射角等于临界角时的光路图，并标明临界角。

\item
当透明介质处在真空中时，根据临界角的定义导出透明介质的折射率 $n$ 与临界角 $\theta_{c}$ 的关系式。
\end{enumerate}

%% answer: （1）光从光密介质射到光疏介质中，折射角为 90° 时的入射角叫做临界角。（2）如图，$\theta_{c}$ 为临界角。（3）$\sin\theta_{c} = \frac{1}{n}$
\jdanswer{
\begin{enumerate}
	\item
	光从光密介质射到光疏介质中，折射角为 $90^\circ$ 时的入射角叫做临界角。

	\item
	光路图如图所示，$\theta_{c}$ 为临界角。

	\item
	$\sin\theta_{c} = \frac{1}{n}$
\end{enumerate}
}

%% memo:
\memoanswer{
（1）光从光密介质射到光疏介质中，折射角为 $90^\circ$ 时的入射角叫做临界角。

（2）光路图如图所示，$\theta_{c}$ 为临界角。

\onepicture[w=4.5cm]{figs/20_an.png}

（3）用 $n$ 表示透明介质的折射率，$\theta_{c}$ 表示临界角。由折射定律
$n\sin\theta_{c} = \sin 90^\circ$，
得 $\sin\theta_{c} = \frac{1}{n}$。
}

\item
%% number: 21
%% paperName: 2000年北京春考第 21 题 12 分
%% typeId: 6
%% type: 计算题
%% chapter: 电路
%% point: 分压与分流
%% method: 无
%% score: 12
%% degree: 550
%% duplicateId: 0
%% body:
AB 两地间铺有通讯电缆，长为 $L$，它是由两条并在一起彼此绝缘的均匀导线组成的，通常称为双线电缆。在一次事故中经检查断定是电缆上某处的绝缘保护层损坏，导致两导线之间漏电，相当于该处电缆的两导线之间接了一个电阻。检查人员经过下面的测量可以确定损坏处的位置：

\begin{enumerate}
\item
令 B 端的双线断开，在 A 处测出双线两端间的电阻 $R_{A}$；

\item
令 A 端的双线断开，在 B 处测出双线两端的电阻 $R_{B}$；

\item
在 A 端的双线间加一已知电压 $U_{A}$，在 B 端用内阻很大的电压表测出两线间的电压 $U_{B}$。
\end{enumerate}
试由以上测量结果确定损坏处的位置。

%% answer: $x = \frac{R_{A}(U_{A} - U_{B})}{(R_{A} + R_{B})U_{A} - 2R_{A}U_{B}}L$
\jdanswer{
$x = \frac{R_{A}(U_{A} - U_{B})}{(R_{A} + R_{B})U_{A} - 2R_{A}U_{B}}L$
}

%% memo:
\memoanswer{
设双线电缆每单位长的电阻为 $r$，漏电处电阻为 $R$，漏电处距 A 端为 $x$，则
$R_{A} = rx + R$，$R_{B} = r(L - x) + R$。
由欧姆定律可知
$\frac{R}{rx + R} = \frac{U_{B}}{U_{A}}$。
解得
$x = \frac{R_{A}(U_{A} - U_{B})}{(R_{A} + R_{B})U_{A} - 2R_{A}U_{B}}L$。
}

\item
%% number: 22
%% paperName: 2000年北京春考第 22 题 12 分
%% typeId: 6
%% type: 计算题
%% chapter: 气体性质
%% point: 气体连接体
%% method: 无
%% score: 12
%% degree: 350
%% duplicateId: 0
%% body:
如图所示，在固定的气缸 A 和 B 中分别用活塞封闭有一定质量的理想气体，活塞面积之比 $S_{A}:S_{B} = 1:2$。两活塞以穿过 B 的底部的刚性细杆相连，可沿水平方向无摩擦滑动，两个气缸都不漏气，初始时 A、B 中气体的体积皆为 $V_{0}$，温度皆为 $T_{0} = 300\UK$。A 中气体压强 $p_{A} = 1.5p_{0}$，$p_{0}$ 是气缸外的大气压强。现对 A 加热，使其中气体的压强升到 $p_{A}' = 2p_{0}$，同时保持 B 中气体的温度不变，求此时 A 中气体温度 $T_{A}'$。
\onepicture[w=8cm, align=right]{figs/22.png}

%% answer: 500 K
\jdanswer{
$500\UK$
}

%% memo:
\memoanswer{
活塞平衡时，有
$p_{A}S_{A} + p_{B}S_{B} = p_{0}(S_{A} + S_{B})$，
$p_{A}'S_{A} + p_{B}'S_{B} = p_{0}(S_{A} + S_{B})$。
已知 $S_{B} = 2S_{A}$，由初始条件可得 $p_{B} = 0.75p_{0}$。
B 中气体初、末态温度相等，设末态体积为 $V_{B}$，则有
$p_{B}'V_{B} = p_{B}V_{0}$。
由末态活塞平衡可得 $p_{B}' = 0.5p_{0}$，故 $V_{B} = 1.5V_{0}$。
设 A 中气体末态体积为 $V_{A}$，因为两活塞移动的距离相等，故有
$\frac{V_{A} - V_{0}}{S_{A}} = \frac{V_{B} - V_{0}}{S_{B}}$，
解得 $V_{A} = 1.25V_{0}$。
由气态方程
$\frac{p_{A}'V_{A}}{T_{A}'} = \frac{p_{A}V_{0}}{T_{A}}$，
解得
$T_{A}' = \frac{p_{A}'V_{A}}{p_{A}V_{0}}T_{A} = 500\UK$。
}

\item
%% number: 23
%% paperName: 2000年北京春考第 23 题 13 分
%% typeId: 6
%% type: 计算题
%% chapter: 运动和力
%% point: 动量和能量
%% method: 无
%% score: 13
%% degree: 450
%% duplicateId: 0
%% body:
云室处在磁感应强度为 $B$ 的匀强磁场中，一静止的质量为 $M$ 的原子核在云室中发生一次 $\alpha$ 衰变，$\alpha$ 粒子的质量为 $m$，电量为 $q$，其运动轨迹在与磁场垂直的平面内。现测得 $\alpha$ 粒子运动的轨道半径 $R$，试求在衰变过程中的质量亏损。（注：涉及动量问题时，亏损的质量可忽略不计）

%% answer: $\Delta m = \frac{Mq^{2}B^{2}R^{2}}{2c^{2}m(M - m)}$
\jdanswer{
$\Delta m = \frac{Mq^{2}B^{2}R^{2}}{2c^{2}m(M - m)}$
}

%% memo:
\memoanswer{
令 $v$ 表示 $\alpha$ 粒子的速度，由洛伦兹力和牛顿定律可得
$qvB = m\frac{v^{2}}{R}$。
令 $v'$ 表示 $\alpha$ 衰变后剩余核的速度，考虑衰变过程中系统动量守恒，因为亏损质量很小，可不予考虑，由动量守恒可知
$(M - m)v' = mv$。
在衰变过程中，$\alpha$ 粒子和剩余核的动能来自于亏损质量，即
$\Delta m \cdot c^{2} = \frac{1}{2}(M - m)v'^{2} + \frac{1}{2}mv^{2}$。
解得
$\Delta m = \frac{Mq^{2}B^{2}R^{2}}{2c^{2}m(M - m)}$。
}

\item
%% number: 24
%% paperName: 2000年北京春考第 24 题 15 分
%% typeId: 6
%% type: 计算题
%% chapter: 运动和力
%% point: 动量和能量
%% method: 无
%% score: 15
%% degree: 350
%% duplicateId: 0
%% body:
相隔一定距离的 A、B 两球，质量相等，假定它们之间存在恒定的斥力作用。原来两球被按住，处在静止状态。现突然松开两球，同时给 A 球以速度 $v_{0}$，使之沿两球连线射向 B 球，B 球初速为零。若两球间的距离从最小值（两球未接触）到刚恢复到原始值所经历的时间为 $t_{0}$。求 B 球在斥力作用下的加速度。

%% answer: $a = \frac{v_{0}}{2t_{0}}$
\jdanswer{
$a = \frac{v_{0}}{2t_{0}}$
}

%% memo:
\memoanswer{
以 $m$ 表示每个球的质量，$F$ 表示恒定斥力，$l$ 表示两球间的原始距离。松开后，A 球做初速为 $v_{0}$ 的匀减速运动，B 球做初速为零的匀加速运动。设在两球间的距离由 $l$ 变小到恢复到 $l$ 的过程中，A 球的路程为 $l_{1}$，B 球的路程为 $l_{2}$；刚恢复到原始长度时，A 球的速度为 $v_{1}$，B 球的速度为 $v_{2}$。由动量守恒定律，有
$mv_{0} = mv_{1} + mv_{2}$。
由功能关系，得
$Fl_{1} = \frac{1}{2}mv_{0}^{2} - \frac{1}{2}mv_{1}^{2}$，$Fl_{2} = \frac{1}{2}mv_{2}^{2}$。
由于初态和末态两球之间的距离相等，故有 $l_{1} = l_{2}$，由以上解得 $v_{2} = v_{0}$。
当两球的速度相等时，距离最小，设此时球的速度为 $u$，则由动量守恒定律，得
$mv_{0} = 2mu$。
设 $a$ 为 B 球的加速度，则有
$v_{2} = u + at_{0}$，
得
$a = \frac{v_{0}}{2t_{0}}$。
}

\end{enumerate}

\end{document}
